% PlanetMath proposal draft for arXiv:2602.08138
% Source corpora: arXiv math.CT eprints, storage/mo-processed-gpu, storage/math-processed-gpu
\section*{Gamified Kat\v{e}tov order and the Effective Topos}
\subsection*{Source}
\begin{itemize}
  \item arXiv:2602.08138, Takayuki Kihara \& Ming Ng, ``The Game-Theoretic Katětov Order and Idealised Effective Subtoposes'' (2026-02-12)
  \item Cross-reference MO 84393 on understanding the Effective Topos and its Lawvere-Tierney topologies
  \item Cross-reference Math.SE 4442510 on interpreting statements external to the Effective Topos
\end{itemize}
\subsection*{Synopsis}
The paper solves a longstanding problem: describing the $\leq_{\mathrm{LT}}$ lattice of Lawvere-Tierney topologies in the Effective Topos $\Eff$. The key idea is a \emph{gamified Katětov order} on filters of $\omega$, where players witness largeness relations. This modified order collapses many classical distinctions (e.g., all MAD families become equivalent) yet retains enough structure to produce infinite ascending chains. The authors show that computable versions of this order are isomorphic to the $\leq_{\mathrm{LT}}$ order, derive a new degree-spectrum invariant $\mathcal{D}_{\mathrm{T}}(\mathcal{F})$, and prove that for any $\Delta^1_1$ filter, the invariant recovers the hyperarithmetic degrees. The result exposes a deep combinatorial core underlying Effective-topos subtoposes.
\subsection*{MathOverflow cues}
\begin{description}
  \item[MO 84393] requests intuition for the Effective Topos; the gamified Katětov correspondence clarifies how its LT topologies mirror concrete filter orders on $\omega$.
\end{description}
\subsection*{Math.SE anchors}
\begin{description}
  \item[Math.SE 4442510] asks how to externalize Effective Topos statements; the paper’s degree-spectrum invariant provides an external characterization of subtoposes definable by $\Delta^1_1$ filters.
\end{description}
\subsection*{Encyclopedia outline}
\begin{enumerate}
  \item \textbf{Classical Katětov vs. gamified version.} Define the new game-based order, compare it to Rudin-Keisler/Tukey orders, and explain collapses/strict chains.
  \item \textbf{Lawvere-Tierney topologies in $\Eff$.} Recall LT topologies, MO 84393’s intuition, and state the isomorphism between $\leq_{\mathrm{LT}}$ and the computable gamified order.
  \item \textbf{Degree-spectrum invariant $\mathcal{D}_{\mathrm{T}}(\mathcal{F})$.} Describe its definition and show how it produces proper initial segments of the Turing degrees.
  \item \textbf{Applications to $\Delta^1_1$ filters.} Prove that hyperarithmetic degrees arise precisely from such filters, addressing Math.SE 4442510’s need for external viewpoints.
  \item \textbf{Implications for computability and topos theory.} Discuss how this bridges combinatorial filter theory with categorical structure.
\end{enumerate}
\subsection*{Action items}
\begin{itemize}
  \item Present the game rules and main lemmas relating gamified Katětov order to LT topologies.
  \item Elaborate on the construction of $\mathcal{D}_{\mathrm{T}}(\mathcal{F})$, including worked examples.
  \item Connect results to known hierarchies (Rudin-Keisler, Tukey) and explain collapses vs. strict chains.
  \item Add cross-links to PlanetMath entries on filters, Lawvere-Tierney topologies, and the Effective Topos.
\end{itemize}
